\section{Endpoint Reconstruction}
\label{sec:endpoint-reconstruction}

After each refined local piece satisfies half-space Stokes, the proof still has
to reconstruct a global represented equality.  This section explains why the
endpoint layer is a mathematical component of the formalization rather than a
renaming of local terms.

\subsection{Nested finite sums}

The selected boundary refinement is indexed in two stages.  There is a finite
set of boundary chart indices \(i\), and for each \(i\) a finite set of active
refined pieces \(q\).  The local theorem gives
\[
  \code{localBulkTerm}(i,q)=\code{localBoundaryTerm}(i,q)
\]
for every active pair.  Summing gives
\[
  \sum_i\sum_{q\in Q_i}\code{localBulkTerm}(i,q)
  =
  \sum_i\sum_{q\in Q_i}\code{localBoundaryTerm}(i,q).
\]

In Lean, this is not merely an algebraic observation.  The finite sets
\(Q_i\) must be definitionally or propositionally connected with the selected
half-space cover, the smooth refinement, and the local Stokes data.  The
endpoint structure is designed to keep those indices aligned.

The guiding invariant is that the same active pair \((i,q)\) performs three
roles at once: it chooses the half-space box, chooses the refined coefficient,
and appears in the endpoint sum.
If any of these are separated, the proof acquires adapter fields whose only
purpose is to prove that two finite descriptions are the same.  The localized
endpoint layer avoids this by generating the endpoint from the refined pieces
that already carry the local Stokes data.

\begin{lemma}[Finite endpoint summation]
\label{lem:endpoint-summation}
Given a localized refined endpoint with finite boundary pieces \(Q_i\), if
\[
  B_{i,q}=C_{i,q}
\]
for every active refined piece \(q\in Q_i\), then the generated represented
endpoints satisfy
\[
  \sum_i\sum_{q\in Q_i} B_{i,q}
  =
  \sum_i\sum_{q\in Q_i} C_{i,q}.
\]
The Lean endpoint layer instantiates \(B_{i,q}\) as
\code{localBulkTerm i q} and \(C_{i,q}\) as
\code{localBoundaryTerm i q}.
\end{lemma}

The lemma looks elementary, but its Lean version is valuable because it fixes
the finite index family once and for all.  The same membership proof
\(q\in Q_i\) is used to access the box corners, the coefficient, the local
form, the zero-support theorem, and the local Stokes fields.  This prevents
the proof from silently changing finite sums when moving between local and
global layers.

The proof itself has three algebraic steps.  First, it rewrites the generated
bulk endpoint as the nested sum of the local bulk terms.  Second, it rewrites
each local bulk term using the per-box Stokes equality.  Third, it folds the
resulting nested sum back into the generated boundary endpoint.  In a paper
proof these are often collapsed into the phrase ``summing over the cover.''
In Lean they are separate because every sum carries its index type and
membership conditions.

\subsection{Localized refined partitions}

The core endpoint structure is
\[
\code{CoverIndexedBoundaryLocalizedRefinedPartition}.
\]
It stores the refined boundary pieces and the generated local bulk and
boundary terms.  Its definitions
\[
\code{generatedRepresentedBulkIntegral}
\]
and
\[
\code{generatedRepresentedBoundaryIntegral}
\]
are the represented endpoints at this layer.  The theorem
\[
\code{representedStokes_of_localStokes}
\]
turns a finite family of local equalities into equality of these generated
endpoints.

The zero-localized route refines this through
\[
\code{CoverIndexedBoundaryLocalizedRefinedPartition.ZeroLocalStokesData}.
\]
This record stores the specific local data needed to apply the half-space
local theorem to zero-localized representatives.  It includes source and
target charts, local forms, box corners, zero-support information, and the
local terms whose equality will be summed.

The localized refined partition is intentionally smaller than the older global
endpoint certificates.  It does not attempt to encode every possible future
native integral comparison.  It records the data needed to state and prove the
represented finite-sum equality.  This makes it a stable target: future
comparison theorems can map from this generated endpoint to native integrals
without changing the local Stokes proof.

This design also prevents the represented theorem from depending on a future
choice of native manifold integration API.  The endpoint structure fixes the
finite coordinate equality now; a later native bridge can prove that this
finite coordinate quantity agrees with a chosen global integral definition.

\subsection{From zero local Stokes data to endpoints}

The theorem
\[
\code{ZeroLocalStokesData.localStokes_of_interiorFields}
\]
builds the per-box local Stokes equality from interior fields and zero support.
Then
\[
\code{boundaryHalfSpaceBulkSum_eq_trueBoundarySum_of_zeroLocalStokesData}
\]
proves the nested finite-sum equality.  Finally,
\[
\code{representedStokes_of_zeroLocalStokesData}
\]
reconstructs equality of the generated represented endpoints.

\begin{theorem}[Zero-localized endpoint reconstruction]
For the localized refined partition generated by the intrinsic route, the
combination of zero support and half-space interior fields implies equality of
the generated represented bulk and boundary endpoints.  In Lean this is the
role of:
\begin{lstlisting}[language=Lean]
CoverIndexedBoundaryLocalizedRefinedPartition.
  representedStokes_of_zeroLocalStokesData
\end{lstlisting}
\end{theorem}

The first-principles route uses the constructor
\[
\code{hasIntrinsicRoute_of_selectedSmoothRefinement}
\]
to assemble the endpoint, zero-localized local data, boundary piece equality,
zero support, and interior fields into an intrinsic certificate.  Once the
canonical selected smooth refinement supplies the last two ingredients, the
intrinsic route has a canonical endpoint equality.

The role of this constructor is precisely to avoid exposing an endpoint
adapter to the final user.  It proves that the endpoint, local Stokes data,
zero-support theorem, and regularity fields are all generated from the same
selected smooth refinement.  Thus the endpoint equality is not imported from
outside the proof; it is assembled from the local Stokes theorem over the
canonical finite route.

\subsection{Canonical endpoints at the top level}

The intrinsic input defines
\[
\code{canonicalRepresentedBulkIntegral}
\quad\text{and}\quad
\code{canonicalRepresentedBoundaryIntegral}
\]
by choosing the route generated from the canonical selected smooth refinement:
\[
\code{hasIntrinsicRoute_intrinsic}.
\]
The first-principles input then projects these endpoints through its generated
intrinsic input.  Thus the top-level definitions are not external names.  They
are the endpoint definitions induced by the canonical finite selected route.

\subsection{Why reconstruction matters}

Endpoint reconstruction is where the represented theorem becomes global.  The
half-space local theorem knows only about one coordinate box.  Zero-localized
support knows only how to kill artificial faces for one refined representative.
Compactness and refinement know how to generate finitely many such pieces.
The endpoint layer proves that all of these local facts assemble into one
equality between the two canonical represented quantities associated with the
input form.

This is also the layer that makes the represented-to-native future work
well-posed.  A future theorem comparing represented endpoints with native
manifold integrals can target
\[
\code{canonicalRepresentedBulkIntegral}
\quad\text{and}\quad
\code{canonicalRepresentedBoundaryIntegral}
\]
directly.  It does not need to rediscover the local proof route.

There is an important conceptual payoff here.  The proof does not define a
global endpoint first and then ask local lemmas to match it.  Instead, the
endpoint is generated from the same local pieces used by the proof.  This
removes a common source of mismatch in large formalizations: the local theorem,
the summation theorem, and the final theorem all quantify over the same finite
piece data.

This is why the word canonical is warranted in the represented theorem.  It
does not mean independent of every possible choice of partition or chart cover;
that would be the native comparison problem.  It means canonical relative to
the generated first-principles route: once the input form is fixed, the route
constructs a selected finite proof object, and the endpoints are the finite
sums determined by that object.

\subsection{What remains for native integrals}

Endpoint reconstruction stops at the finite coordinate equality.  To obtain a
future native theorem, one must prove two comparison results.  The bulk
comparison identifies the represented bulk sum with the chosen definition of
\(\int_M\dd\omega\).  The boundary comparison identifies the represented
lower-face sum with the chosen definition of \(\int_{\partial M}\omega\) under
the global boundary orientation.  These comparisons are logically downstream
from the represented theorem: they should not be mixed into the proof that the
finite coordinate endpoints are equal.
